\chapter{Комбинаторика}
\section{Перестановки}
\subsection{Факториал}
\import{factorial.tex}
\kactlimport{IntPerm.h}
\subsection{Циклы}
Число перестановок, длины циклов которых лежат в S:
\[\sum_{n\ge0}g_S(n)\frac{x^n}{n!}=\exp\left(\sum_{k\in S}\frac{x^k}{k}\right).\]
\subsection{Беспорядки}
Перестановки без неподвижных точек: $D_0=1$, $D_1=0$;
\[D_n=(n-1)(D_{n-1}+D_{n-2})=nD_{n-1}+(-1)^n.\]
При n>=1 это ближайшее целое к n!/e.
\subsection{Лемма Бёрнсайда}
Число орбит множества X под действием группы G:
\[\frac1{|G|}\sum_{g\in G}|X^g|,\]
где $X^g$ — элементы, неподвижные при g.
Для конфигураций длины n с точностью до циклических сдвигов:
\[g(n)=\frac1n\sum_{k=0}^{n-1}f(\gcd(n,k))
=\frac1n\sum_{d\mid n}f(d)\phi(n/d),\]
если f(d) считает допустимые периодические конфигурации периода,
делящего d.
\section{Разбиения и подмножества}
\subsection{Число разбиений}
Число представлений n суммой положительных слагаемых без учёта порядка:
\[p(0)=1,\quad p(n)=\sum_{k\in\mathbb Z\setminus\{0\}}
(-1)^{k+1}p\left(n-\frac{k(3k-1)}2\right).\]
При отрицательном аргументе p=0.
$p(n)\sim\exp(\pi\sqrt{2n/3})/(4n\sqrt3)$.
Значения для n=0..9: 1,1,2,3,5,7,11,15,22,30;
$p(20)=627$, $p(50)\approx2\cdot10^5$, $p(100)\approx2\cdot10^8$.
\subsection{Теорема Лукаса}
Для простого p записать n,m по основанию p:
$n=\sum n_ip^i$, $m=\sum m_ip^i$. Тогда
$\binom nm\equiv\prod_i\binom{n_i}{m_i}\pmod p$.
\subsection{Биномиальные коэффициенты}
\kactlimport{FastCombinatorics.h}
\kactlimport{multinomial.h}
\section{Специальные числа}
\subsection{Числа Бернулли}
ЭПФ: $B(t)=t/(e^t-1)$.
$B_0,B_1,\dots=1,-1/2,1/6,0,-1/30,0,1/42,\dots$.
\[\sum_{i=1}^n i^m=\frac1{m+1}\sum_{k=0}^{m}
\binom{m+1}{k}B_k(n+1)^{m+1-k}\quad(m\ge1).\]
Формула Эйлера — Маклорена для затухающей функции:
\begin{align*}
\sum_{i=m}^{\infty}f(i)&\sim\int_m^\infty f(x)\,dx-
\sum_{k\ge1}\frac{B_k}{k!}f^{(k-1)}(m)\\
&=\int_m^\infty f(x)\,dx+\frac{f(m)}2-\frac{f'(m)}{12}
+\frac{f'''(m)}{720}-\dots
\end{align*}
\subsection{Числа Стирлинга первого рода}
Число перестановок n элементов с k циклами:
\[c(n,k)=c(n-1,k-1)+(n-1)c(n-1,k),\quad c(0,0)=1.\]
$\sum_kc(n,k)x^k=x(x+1)\dots(x+n-1)$.
$c(8,k)$ для k=0..8:
0,5040,13068,13132,6769,1960,322,28,1.
$c(n,2)$ для n=0..9: 0,0,1,3,11,50,274,1764,13068,109584.
\subsection{Эйлеровы числа}
Число перестановок с k спусками ($\pi_j>\pi_{j+1}$):
\begin{align*}
E(n,k)&=(n-k)E(n-1,k-1)+(k+1)E(n-1,k),\\
E(n,0)&=E(n,n-1)=1,\\
E(n,k)&=\sum_{j=0}^{k}(-1)^j\binom{n+1}{j}(k+1-j)^n.
\end{align*}
То же число перестановок имеет k превышений ($\pi_j>j$)
или k+1 слабых превышений ($\pi_j\ge j$).
\subsection{Числа Стирлинга второго рода}
Разбиения n различных элементов на k непустых групп:
\begin{align*}
S(n,k)&=S(n-1,k-1)+kS(n-1,k),\\
S(n,1)&=S(n,n)=1,\quad S(0,0)=1,\\
S(n,k)&=\frac1{k!}\sum_{j=0}^{k}(-1)^{k-j}\binom kjj^n.
\end{align*}
\subsection{Числа Белла}
Все разбиения n различных элементов: $B_n=\sum_kS(n,k)$.
Значения: 1,1,2,5,15,52,203,877,4140,21147,\dots.
Для простого p: $B_{p^m+n}\equiv mB_n+B_{n+1}\pmod p$.
\subsection{Помеченные деревья}
На n вершинах: $n^{n-2}$.
Число способов соединить k заданных деревьев размеров $n_i$
в одно: $n_1\dots n_k\,n^{k-2}$, $n=\sum n_i$.
С заданными степенями $d_i$:
$(n-2)!/\prod_i(d_i-1)!$ при $\sum d_i=2n-2$.
\subsection{Числа Каталана}
\begin{align*}
C_n&=\frac1{n+1}\binom{2n}{n}=\binom{2n}{n}-\binom{2n}{n+1},\\
C_0&=1,\quad C_{n+1}=\frac{2(2n+1)}{n+2}C_n
=\sum_{i=0}^n C_iC_{n-i}.
\end{align*}
Значения: 1,1,2,5,14,42,132,429,1430,4862,16796,\dots.
Считают: правильные скобочные последовательности длины 2n;
пути в квадрате $n\times n$ под диагональю; полные бинарные деревья
с n+1 листьями; упорядоченные корневые деревья на n+1 вершинах;
триангуляции выпуклого (n+2)-угольника; перестановки n элементов
без возрастающей подпоследовательности длины 3.
